%% Chapter 9 — The Reasoning Receipt Chain as Epistemic Infrastructure

\chapter{The Reasoning Receipt Chain as Epistemic Infrastructure}
\label{chap:receipt}

\epigraph{%
  ``Trust, but verify.''
}{--- Russian proverb, adopted by Ronald Reagan; for AI, trust the algorithm, verify the execution}

Chapter~\ref{chap:speed} established that sub-microsecond breed latency makes it possible
to manufacture millions of reasoning conclusions within a single decision budget.  This
capability introduces an accountability obligation that has no precedent in classical expert
systems: when a system produces $3 \times 10^7$ MYCIN consultations per minute, every one
of those consultations must be individually auditable, non-repudiable, and causally traceable
to the exact input that produced it.

This chapter formalises the \emph{reasoning receipt chain} --- the cryptographic
infrastructure that satisfies this obligation.  We prove that BLAKE3 hashing constitutes
a process-binding commitment scheme (Section~\ref{sec:blake3-commitment}), that receipt
composition is functorial over breed morphisms
(Section~\ref{sec:receipt-functor}), and that the receipt chain satisfies four
civilizational trust properties sufficient for regulatory compliance
(Section~\ref{sec:trust-properties}).

%% -----------------------------------------------------------------------
\section{BLAKE3 as a Commitment Scheme}
\label{sec:blake3-commitment}
%% -----------------------------------------------------------------------

\begin{definition}[Reasoning Receipt]
\label{def:reasoning-receipt}
Let $\Breed_i$ be a breed, $A$ an input assertion set, $L = (l_1, \ldots, l_n)$ the
ordered sequence of rules fired during the consultation, and $t$ a Unix nanosecond
timestamp.  The \emph{reasoning receipt} is
\[
  \rho \;=\; \blake(\mathtt{breed\_id}(\Breed_i) \;\|\; A \;\|\; L \;\|\; t),
\]
where $\blake$ denotes the BLAKE3 hash function, $\|$ denotes concatenation after
canonical serialisation (JSON with lexicographically sorted keys), and
$\mathtt{breed\_id}(\Breed_i)$ is the stable identifier registered in the kernel registry.
\end{definition}

The receipt $\rho$ is a 256-bit digest.  Three properties make it suitable as a
commitment scheme for reasoning outputs.

\begin{definition}[Three Receipt Properties]
\label{def:receipt-properties}
A reasoning receipt $\rho$ satisfies:
\begin{enumerate}[label=(\roman*)]
  \item \emph{Binding.}  It is computationally infeasible to find
        $(A', L', t') \neq (A, L, t)$ such that
        $\blake(\mathtt{breed\_id} \| A' \| L' \| t') = \rho$.
        This follows immediately from the collision resistance of BLAKE3 under
        the standard random-oracle assumption.

  \item \emph{Hiding.}  Given only $\rho$, it is computationally infeasible to
        recover $A$, $L$, or $t$.  This follows from the preimage resistance of BLAKE3
        and the unpredictability of $t$ at nanosecond resolution.

  \item \emph{Process-binding.}  The receipt commits not merely to the input--output
        pair $(A, \text{conclusion})$ but to the entire \emph{execution trace} $L$,
        i.e., the ordered sequence of rules fired.  Two consultations that reach the
        same conclusion by different rule paths produce different receipts.
\end{enumerate}
\end{definition}

\paragraph{Causal completeness update.}
The original formulation of Definition~\ref{def:reasoning-receipt} covered the breed
identifier, the rule-firing trace, and the timestamp.  A gap remained: the serialised
input assertion set $A$ was available to the TypeScript consumer but was not independently
hashed at the WASM boundary, creating the possibility that a receipt could be replayed
with a silently modified input.

This gap is closed in wasm4pm v26.6.10.  The function \texttt{cognition\_run()} in
\texttt{crates/wasm4pm-cognition/src/wasm.rs} now computes
\[
  \mathtt{input\_hash} \;=\; \blake(\mathtt{input\_json}),
\]
where \texttt{input\_json} is the raw byte string received at the WASM boundary before
any deserialisation.  The field \texttt{input\_hash} is included in the
\texttt{ContractResult} struct returned to the TypeScript consumer and is validated for
presence by the schema at \texttt{packages/cognition/src/schemas.ts:128}.  The receipt
therefore now reads
\[
  \rho \;=\; \blake(\mathtt{breed\_id} \;\|\; \mathtt{input\_hash} \;\|\; L \;\|\; t),
\]
and the causal completeness gap is closed: any modification to the input, however small,
changes \texttt{input\_hash}, and therefore changes $\rho$, making the tamper-evident
property unconditional rather than dependent on the integrity of the TypeScript layer.

\paragraph{Process-binding vs.\ zero-knowledge proof.}
A natural question is whether a zero-knowledge proof (ZKP) of correct execution would
be preferable to a process-binding receipt.  A ZKP would allow a verifier to confirm that
the computation was performed correctly without learning the inputs.  However, regulatory
compliance frameworks (FDA 21 CFR Part 11, EU AI Act Article 13, ISO 13485) require that
\emph{both} the input record and the execution trace be available to an authorised auditor;
hiding them behind a ZKP would violate the transparency obligations those frameworks impose.
Process-binding commitment is therefore the correct primitive for regulated AI: it provides
computational integrity guarantees (no one can forge a receipt for a computation they did
not perform) while preserving full auditability for authorised parties.

%% -----------------------------------------------------------------------
\section{The Reasoning Receipt Functor}
\label{sec:receipt-functor}
%% -----------------------------------------------------------------------

We now show that receipt generation is not an ad-hoc operation but a natural transformation
in the categorical framework established in Chapter~\ref{chap:breed}.

\begin{definition}[Receipt Functor]
\label{def:receipt-functor}
Define the \emph{receipt functor}
\[
  R_{\Breed} \;:\; \mathbf{Breed} \;\longrightarrow\; \mathbf{Receipt}
\]
as follows.  On objects: $R_{\Breed}(\Breed_i) = \rho_i$, the receipt of the terminal
consultation of $\Breed_i$.  On morphisms: if $f : \Breed_i \to \Breed_j$ is a breed
morphism (a natural transformation of the underlying inference functors), then
$R_{\Breed}(f) : \rho_i \to \rho_j$ is the canonical hash-linking map
$\rho_j = \blake(\rho_i \| \mathtt{breed\_id}(\Breed_j) \| \ldots)$, i.e., the receipt
of the successor breed includes the receipt of the predecessor as a committed prefix.
\end{definition}

The Kleisli-composition interpretation makes the chain structure transparent.  Let
$\mathbf{Receipt}$ be the category whose objects are 256-bit digests and whose morphisms
are hash-linking maps.  Then $R_{\Breed}$ maps a sequential pipeline of breed invocations
$\Breed_1 \to \Breed_2 \to \cdots \to \Breed_n$ to a receipt chain
$\rho_1 \to \rho_2 \to \cdots \to \rho_n$ in which each receipt is a committed function
of all predecessors.  This is precisely the Kleisli composition in the writer monad over
the free monoid of BLAKE3 digests: composing two breed invocations appends their receipts,
and the identity morphism maps to the empty receipt chain.

\begin{theorem}[Civilizational Trust Properties]
\label{thm:civilizational-trust}
The receipt functor $R_{\Breed}$ satisfies the following four properties for any pipeline
of breed invocations operating under the wasm4pm kernel:

\begin{enumerate}[label=(\roman*)]
  \item \emph{Non-forgeability.}
  \item \emph{Non-repudiation.}
  \item \emph{Third-party auditability.}
  \item \emph{Causal completeness.}
\end{enumerate}
\end{theorem}

\begin{proof}
We prove each property in turn.

\paragraph{(i) Non-forgeability.}
Suppose an adversary wishes to present a receipt $\rho'$ claiming that a particular breed
$\Breed_i$ was invoked on input $A'$ and fired rule sequence $L'$, without having actually
performed that invocation.  To produce $\rho'$, the adversary must find a preimage of the
BLAKE3 output
$\rho' = \blake(\mathtt{breed\_id}(\Breed_i) \| \mathtt{input\_hash'} \| L' \| t')$.
Under the standard collision-resistance and preimage-resistance assumptions for BLAKE3
(which is based on the ARX permutation and has received no known attack reducing its
security below 128 bits), this is computationally infeasible.  Furthermore, the
\texttt{input\_hash} field is computed at the WASM boundary before any TypeScript code
can modify the input, so the adversary cannot substitute a different input after the hash
is committed.  The chain structure amplifies this: forging any receipt in the middle of a
chain would require forging all subsequent receipts simultaneously, multiplying the
computational difficulty by the chain length.

\paragraph{(ii) Non-repudiation.}
Once a receipt $\rho$ is emitted by \texttt{cognition\_run()}, the issuer (the wasm4pm
runtime) cannot plausibly deny that the invocation occurred.  The receipt is saved to
\texttt{.wasm4pm/receipts/latest.json} and, in compliant deployments, to an append-only
audit log.  Because the receipt commits to \texttt{breed\_id}, \texttt{input\_hash}, the
rule-firing trace $L$, and the timestamp $t$, any claim by the issuer that the invocation
did not occur, or occurred with different parameters, can be immediately falsified by
presenting the receipt and replaying the computation: the same inputs must produce the
same receipt (determinism is a merge gate in wasm4pm; see \texttt{CLAUDE.md}).  The
non-repudiation property therefore rests on the combination of BLAKE3 binding and the
determinism invariant.

\paragraph{(iii) Third-party auditability.}
A third-party auditor possessing only the receipt chain $(\rho_1, \ldots, \rho_n)$ and
the original input-output pairs can independently verify the chain's integrity by
recomputing each $\rho_i$ from the committed data.  The auditor requires no access to
the runtime environment, no secret keys, and no privileged channels.  This is the
\emph{process-binding} property of Definition~\ref{def:receipt-properties}(iii): the
receipt encodes not just the conclusion but the execution path, so the auditor can
distinguish between two consultations that reached the same conclusion by different rule
paths --- a distinction that zero-knowledge proofs could not provide without disclosing
the inputs.  Third-party auditability is therefore a strictly stronger property than
output correctness: it certifies \emph{how} the conclusion was reached, not merely
\emph{that} the conclusion is consistent with the inputs.

\paragraph{(iv) Causal completeness.}
Prior to v26.6.10, a receipt committed to the breed identifier, the rule-firing trace,
and the timestamp, but the input assertion set $A$ was hashed only in the TypeScript
consumer layer.  This created a causal gap: a receipt could in principle be replayed
against a modified input if the TypeScript layer were compromised, without changing the
receipt value.  The causal completeness gap is now closed by the addition of
\texttt{input\_hash} $= \blake(\mathtt{input\_json})$ computed in
\texttt{crates/wasm4pm-cognition/src/wasm.rs} before deserialisation.  With this
addition, the receipt is a committed function of the complete causal closure of the
computation: the breed, the input, the execution trace, and the time.  Any modification
to any of these four elements changes the receipt.  Causal completeness means that the
receipt chain is a \emph{sufficient statistic} for the computation: given the chain,
an auditor can reconstruct the entire causal history without any additional evidence.
This is the epistemic property that justifies calling the receipt chain
\emph{infrastructure} rather than merely a log: it is the substrate on which
the trust properties of (i)--(iii) rest.
\end{proof}

%% -----------------------------------------------------------------------
\section{Summary}
\label{sec:receipt-summary}
%% -----------------------------------------------------------------------

This chapter proved that the wasm4pm reasoning receipt chain is not an engineering
convenience but an epistemic infrastructure satisfying four civilizational trust properties:
non-forgeability, non-repudiation, third-party auditability, and --- following the v26.6.10
\texttt{input\_hash} addition at the WASM boundary --- causal completeness.  Together,
Chapters~\ref{chap:speed} and~\ref{chap:receipt} establish the two sides of the
accountability equation: the speed phase transitions of Chapter~\ref{chap:speed} make
ensemble reasoning feasible; the receipt functor of this chapter makes every conclusion in
that ensemble auditable.  Chapter~\ref{chap:synthesis} draws these threads together into
the Periodic Table of Reason, showing how the categorical, process-mining, and
cryptographic constraints jointly determine the shape of trustworthy AI reasoning systems.
